\chapter{Agent-to-Agent Communication (A2A)}
\label{sec:a2a}

As large language models evolve from isolated assistants into collaborative networks of specialized agents, the question of \emph{how agents talk to each other} becomes as important as how they reason internally. This section covers the protocols, patterns, and engineering practices that enable multi-agent systems to coordinate, delegate, and collectively solve problems that no single agent could handle alone.

\section{Motivation: Why Agents Must Communicate}
\label{sec:a2a:motivation}

\begin{intuitionbox}[The Specialization Imperative]
A single generalist agent faces a fundamental tension: breadth of knowledge versus depth of capability. Real-world tasks---legal document review, multi-step scientific research, enterprise software development---demand both. Agent-to-agent communication resolves this tension by allowing a \emph{network} of specialists to collaborate, each contributing its strengths while delegating weaknesses.
\end{intuitionbox}

Several forces drive the need for structured inter-agent communication:

\paragraph{Cognitive Load and Context Limits.}
\label{cognitive-load-and-context-limits.}

Every LLM operates within a finite context window. Complex workflows---spanning hundreds of documents, tool calls, and reasoning steps---quickly exceed what a single agent can hold in memory. By decomposing tasks across agents, each agent operates within a manageable context, and the orchestrating agent maintains only high-level state.

\paragraph{Specialization and Expertise.}
\label{specialization-and-expertise.}

Different agents may be fine-tuned, prompted, or tool-equipped for specific domains: a \texttt{CodeAgent} with access to compilers and test runners, a \texttt{LegalAgent} with access to case-law databases, a \texttt{DataAgent} with statistical libraries. Routing subtasks to the right specialist improves both quality and efficiency.

\paragraph{Parallelism and Throughput.}
\label{parallelism-and-throughput.}

Independent subtasks can be dispatched to multiple agents simultaneously. A research orchestrator might fan out literature searches across five specialized agents in parallel, then synthesize their results---dramatically reducing wall-clock time.

\paragraph{Fault Isolation and Resilience.}
\label{fault-isolation-and-resilience.}

When one agent fails, a well-designed multi-agent system can retry with a different agent, fall back to a simpler approach, or escalate to a human---without collapsing the entire workflow.

\paragraph{Delegation and Handoff.}
\label{delegation-and-handoff.}

Long-running tasks may need to be handed off between agents as context shifts. An initial \texttt{PlannerAgent} decomposes a goal, hands subtasks to \texttt{ExecutorAgents}, and a final \texttt{ReviewerAgent} validates outputs---each agent receiving exactly the context it needs.

\begin{keybox}[Core Requirements for A2A Communication]
\begin{enumerate}
  \item \textbf{Discoverability}: Agents must be able to find other agents and understand their capabilities.
  \item \textbf{Interoperability}: Agents built by different teams or vendors must speak a common protocol.
  \item \textbf{Asynchrony}: Long-running tasks must not block the caller; results arrive via callbacks or polling.
  \item \textbf{Security}: Agents must authenticate each other and enforce authorization boundaries.
  \item \textbf{Observability}: Every message exchange must be traceable for debugging and auditing.
\end{enumerate}
\end{keybox}

\section{The Google A2A Protocol}
\label{sec:a2a:google}

In April 2025, Google (with contributions from over 50 technology partners) released the \textbf{Agent-to-Agent (A2A) Protocol}~\cite{google-a2a-2025}, an open specification for interoperable communication between AI agents. The protocol was subsequently donated to the \textbf{Linux Foundation} and has grown to over 150 supporting organizations as of 2026. A2A is designed around a set of core principles that distinguish it from earlier ad-hoc approaches.

\subsection{Design Philosophy}
\label{design-philosophy}

The A2A specification articulates five guiding principles (adapted from the official spec~\cite{google-a2a-2025}, §1.2):

\begin{keybox}[A2A Design Principles]
Opaque execution

Calling agents never inspect the internals of a remote agent---they interact solely through the declared interface. Whether the target is GPT-4, Gemini, or a rule-based system is irrelevant to the protocol, enabling genuinely heterogeneous agent ecosystems.

Enterprise readiness

Authentication (OAuth 2.0, API keys, JWT), audit logging, and regulatory compliance are not afterthoughts---they are integrated at the protocol level from the outset.

Modality agnosticism

A single message may combine text, binary files, and structured JSON payloads, accommodating agents that operate on images, audio, code, or documents without protocol extensions.

Simplicity via existing standards

Rather than inventing new transports, A2A reuses HTTP/HTTPS with JSON-RPC~2.0 messages, Server-Sent Events (SSE) for streaming, and gRPC as an alternative binding---technologies that every infrastructure team already operates.

Async-first task model

Long-running operations are the norm, not the exception. Push notifications and polling are both first-class mechanisms, so callers never need to hold open a connection for hours.
\end{keybox}

\subsection{Agent Cards}
\label{agent-cards}

The foundation of A2A discoverability is the \textbf{Agent Card}---a machine-readable JSON manifest hosted at a well-known endpoint (\texttt{/.well-known/agent.json}). It advertises what the agent can do, how to authenticate, and where to send tasks---analogous to an OpenAPI spec but for autonomous agents rather than REST endpoints.

\begin{examplebox}[Agent Card Structure]
\begin{lstlisting}[style=pythonstyle]
# Agent Card served at https://agent.example.com/.well-known/agent.json
agent_card = {
    "name": "DataAnalysisAgent",
    "description": "Analyzes structured datasets, produces statistical summaries, "
                   "generates visualizations, and answers data questions.",
    "url": "https://agent.example.com/a2a",
    "version": "1.2.0",
    "capabilities": {
        "streaming": True,
        "pushNotifications": True,
        "stateTransitionHistory": True
    },
    "authentication": {
        "schemes": ["Bearer", "ApiKey"]
    },
    "skills": [
        {
            "id": "statistical-analysis",
            "name": "Statistical Analysis",
            "description": "Compute descriptive statistics, run hypothesis tests, "
                           "fit regression models on tabular data.",
            "tags": ["statistics", "data", "analysis", "regression"],
            "examples": [
                "What is the correlation between columns A and B?",
                "Run a t-test comparing these two groups.",
                "Fit a linear regression predicting sales from ad spend."
            ],
            "inputModes": ["text", "data"],
            "outputModes": ["text", "data", "file"]
        },
        {
            "id": "visualization",
            "name": "Data Visualization",
            "description": "Generate charts, plots, and dashboards from data.",
            "tags": ["charts", "plots", "visualization", "dashboard"],
            "examples": [
                "Create a bar chart of monthly revenue.",
                "Plot the distribution of customer ages."
            ],
            "inputModes": ["text", "data"],
            "outputModes": ["file", "text"]
        }
    ],
    "defaultInputModes": ["text"],
    "defaultOutputModes": ["text"]
}
\end{lstlisting}
\end{examplebox}

Agent Cards enable \emph{capability-based routing}: an orchestrator agent can fetch cards from a registry, semantically match a subtask to the most appropriate agent, and dispatch accordingly---all without hardcoded routing logic.

\subsection{Task Lifecycle}
\label{task-lifecycle}

A2A models all work as \textbf{Tasks}. A task progresses through a well-defined state machine:

\texttt{submitted}

The client has sent the task; the server has acknowledged receipt.

\texttt{working}

The agent is actively processing. The client may poll or await SSE events.

\texttt{input-required}

The agent needs additional information from the user or calling agent before it can proceed (e.g., a clarifying question, a missing credential).

\texttt{completed}

The task finished successfully; results are available in the response.

\texttt{failed}

An unrecoverable error occurred; an error message explains the cause.

\texttt{rejected}

The agent declined the task (e.g., outside its capabilities or unauthorized). Added in A2A v1.0.

\texttt{canceled}

The task was aborted, either by the client or by the server.

\subsection{Streaming via Server-Sent Events}
\label{streaming-via-server-sent-events}

For tasks that produce incremental output (e.g., a long report being written, a code file being generated), A2A uses \textbf{Server-Sent Events (SSE)}. The client opens a persistent HTTP connection and receives a stream of JSON events:

\begin{examplebox}[SSE Event Stream Example]
\begin{lstlisting}[style=pythonstyle]
# Each SSE event carries a TaskStatusUpdateEvent or TaskArtifactUpdateEvent
# Example stream for a "write a research report" task:

# Event 1: status update
data: {
  "id": "task-abc123",
  "status": {"state": "working"},
  "final": false
}

# Event 2: partial artifact (streaming text)
data: {
  "id": "task-abc123",
  "artifact": {
    "parts": [{"type": "text", "text": "## Introduction\n\nRecent advances in..."}],
    "index": 0,
    "append": false,
    "lastChunk": false
  },
  "final": false
}

# Event 3: more text appended
data: {
  "id": "task-abc123",
  "artifact": {
    "parts": [{"type": "text", "text": " reinforcement learning have shown..."}],
    "index": 0,
    "append": true,   # append to existing artifact
    "lastChunk": false
  },
  "final": false
}

# Final event: task complete
data: {
  "id": "task-abc123",
  "status": {"state": "completed"},
  "final": true
}
\end{lstlisting}
\end{examplebox}

\subsection{Push Notifications for Long-Running Tasks}
\label{push-notifications-for-long-running-tasks}

When a task may take minutes or hours, maintaining an open SSE connection is impractical. A2A supports \textbf{push notifications}: the client registers a webhook URL, and the server POSTs status updates as the task progresses.

\begin{lstlisting}[style=pythonstyle]
# Client registers a push notification endpoint when submitting the task
task_request = {
    "id": "task-xyz789",
    "message": {
        "role": "user",
        "parts": [{"type": "text", "text": "Analyze Q3 sales data and produce a report."}]
    },
    "pushNotification": {
        "url": "https://my-orchestrator.example.com/webhooks/a2a",
        "token": "secret-hmac-token-for-verification",
        "authentication": {
            "schemes": ["Bearer"],
            "credentials": "eyJhbGciOiJIUzI1NiJ9..."
        }
    }
}
# The server will POST TaskStatusUpdateEvent objects to the webhook URL
# as the task transitions through states.
\end{lstlisting}

\subsection{Message Format}
\label{message-format}

A2A messages consist of a \textbf{role} (\texttt{user} or \texttt{agent}) plus a list of typed \textbf{parts} (text, file, or structured data). The full message schema, multi-modal examples, and context-passing guidelines are covered in Section~\ref{sec:a2a:messages}.

\subsection{Authentication and Authorization}
\label{authentication-and-authorization}

A2A supports multiple authentication schemes, declared in the Agent Card and enforced per-request:

\begin{itemize}
  \item \textbf{Bearer tokens (JWT/OAuth 2.0)}: Standard for enterprise deployments; tokens carry scopes that limit what the calling agent is permitted to request.
  \item \textbf{API keys}: Simpler scheme for internal or trusted environments.
  \item \textbf{Mutual TLS (mTLS)}: Certificate-based authentication for high-security deployments.
  \item \textbf{OpenID Connect}: Federated identity, enabling cross-organization agent communication.
\end{itemize}

\begin{warningbox}[Authorization Scope Enforcement]
An agent receiving a task must verify not only \emph{who} is calling (authentication) but \emph{what they are allowed to request} (authorization). A \texttt{ReportingAgent} might accept read-only data queries from any authenticated agent, but restrict write operations to agents holding a specific OAuth scope. Failing to enforce this creates privilege escalation vulnerabilities in multi-agent systems.
\end{warningbox}

\section{Communication Patterns}
\label{sec:a2a:patterns}

Multi-agent systems employ a variety of communication patterns depending on the nature of the task, latency requirements, and the number of agents involved.

\subsection{Request-Response}
\label{request-response}

The simplest pattern: Agent A sends a task to Agent B and waits for a complete response. Suitable for short, well-defined subtasks where the result is needed before proceeding.

\subsection{Streaming}
\label{streaming}

Agent A opens an SSE connection; Agent B streams partial results as they are produced. Ideal for long-form generation (reports, code), real-time collaboration, or progressive UI updates.

\begin{examplebox}[Streaming Pattern Use Case]
An orchestrator asks a \texttt{WritingAgent} to draft a 10-page technical document. Rather than waiting 2 minutes for the complete document, the orchestrator streams each section as it is written, allowing a \texttt{ReviewAgent} to begin reviewing early sections while later sections are still being generated---a pipeline that reduces total latency by 40--60\%.
\end{examplebox}

\subsection{Multi-Turn Interaction}
\label{multi-turn-interaction}

Some tasks require iterative refinement. The agent enters \texttt{input-required} state, the orchestrator provides clarification, and the task resumes. This mirrors human collaborative workflows: draft $\to$ feedback $\to$ revision.

\begin{lstlisting}[style=pythonstyle]
# Multi-turn: orchestrator handles input-required state
async def run_multiturn_task(client, initial_message):
    task = await client.send_task(message=initial_message)

    while task.status.state not in ("completed", "failed", "canceled"):
        if task.status.state == "input-required":
            # Agent needs clarification
            clarification_needed = task.status.message
            print(f"Agent asks: {clarification_needed}")

            # Orchestrator generates or forwards a clarifying response
            user_reply = await get_clarification(clarification_needed)

            # Send the reply to continue the task
            task = await client.send_task(
                task_id=task.id,
                message={"role": "user",
                         "parts": [{"type": "text", "text": user_reply}]}
            )
        else:
            # Still working --- poll after a delay
            await asyncio.sleep(2)
            task = await client.get_task(task.id)

    return task
\end{lstlisting}

\subsection{Broadcast}
\label{broadcast}

An orchestrator sends the same message to multiple agents simultaneously---useful for announcements, distributing shared context, or triggering parallel independent workflows.

\subsection{Publish-Subscribe (Pub-Sub)}
\label{publish-subscribe-pub-sub}

Agents subscribe to event channels (e.g., \texttt{new-document-uploaded}, \texttt{model-retrained}). When an event fires, all subscribed agents are notified. This decouples producers from consumers and enables reactive, event-driven architectures.

\subsection{Negotiation}
\label{negotiation}

Two agents exchange proposals and counter-proposals to reach agreement on a plan, resource allocation, or approach. Common in multi-agent planning systems where agents have different objectives or constraints.

\begin{examplebox}[Negotiation Pattern]
A \texttt{PlannerAgent} proposes a 5-step research plan. A \texttt{ResourceAgent} responds that Step 3 (running a large simulation) would exceed the compute budget. The \texttt{PlannerAgent} counter-proposes a scaled-down simulation. The \texttt{ResourceAgent} approves. The agreed plan is then dispatched to executor agents.
\end{examplebox}

\subsection{Auction-Based Task Allocation}
\label{auction-based-task-allocation}

The orchestrator announces a task with requirements; candidate agents submit bids (estimated completion time, confidence, cost); the orchestrator awards the task to the winning bidder. This enables dynamic, market-based load balancing across a pool of agents.

\begin{table}[ht!]
\centering
\caption{Summary of A2A communication patterns.}
\begin{tabular}{@{}lp{5cm}p{8cm}@{}}
\toprule
\textbf{Pattern} & \textbf{Latency} & \textbf{Best For} \\
\midrule
Request-Response & Low & Short, well-defined subtasks \\
Streaming & Low (first token) & Long-form generation, real-time UI \\
Multi-Turn & Medium & Ambiguous tasks requiring clarification \\
Broadcast & Low & Shared context distribution \\
Pub-Sub & Variable & Event-driven reactive workflows \\
Negotiation & Medium--High & Resource-constrained planning \\
Auction & Medium & Dynamic load balancing \\
\bottomrule
\end{tabular}
\end{table}

\section{Agent Discovery and Routing}
\label{sec:a2a:discovery}

Before an agent can communicate with another, it must \emph{find} it. Agent discovery is the process of locating agents that can handle a given task.

\subsection{Agent Registries}
\label{agent-registries}

An \textbf{agent registry} is a directory service that indexes Agent Cards and provides search and lookup APIs. Two deployment models exist:

Centralized Registry

A single authoritative registry (e.g., an enterprise service catalog) indexes all agents. Simple to operate but creates a single point of failure and may not scale to cross-organization deployments.

Federated Registry

Multiple registries, each authoritative for a domain or organization, with cross-registry search protocols. More resilient and privacy-preserving, but requires standardized federation protocols.

\subsection{Capability-Based Routing}
\label{capability-based-routing}

Rather than hardcoding agent URLs, orchestrators perform \textbf{capability-based routing}: they query the registry for agents matching required skills, then select the best match.

\begin{lstlisting}[style=pythonstyle]
class AgentRouter:
    """Routes tasks to agents based on capability matching."""

    def __init__(self, registry_url: str):
        self.registry_url = registry_url
        self._cache: dict[str, list[AgentCard]] = {}

    async def find_agents(self, required_skill: str,
                          tags: list[str] | None = None) -> list[AgentCard]:
        """Query registry for agents with the required skill."""
        params = {"skill": required_skill}
        if tags:
            params["tags"] = ",".join(tags)
        async with httpx.AsyncClient() as client:
            resp = await client.get(f"{self.registry_url}/agents", params=params)
            return [AgentCard(**card) for card in resp.json()["agents"]]

    async def route(self, task_description: str) -> AgentCard:
        """Semantically match a task description to the best available agent."""
        # Embed the task description
        task_embedding = await embed(task_description)

        # Fetch all registered agents
        all_agents = await self.find_agents(required_skill="*")

        # Score each agent by cosine similarity of task to agent description
        scored = []
        for agent in all_agents:
            agent_embedding = await embed(agent.description)
            score = cosine_similarity(task_embedding, agent_embedding)
            scored.append((score, agent))

        # Return the highest-scoring agent
        scored.sort(key=lambda x: x[0], reverse=True)
        return scored[0][1]
\end{lstlisting}

\subsection{Load Balancing Across Equivalent Agents}
\label{load-balancing-across-equivalent-agents}

When multiple agents offer the same capability, the router must distribute load. Common strategies:

\begin{itemize}
  \item \textbf{Round-robin}: Distribute tasks evenly across all available agents.
  \item \textbf{Least-loaded}: Route to the agent with the fewest active tasks (requires health/metrics endpoints).
  \item \textbf{Latency-aware}: Route to the agent with the lowest recent response time.
  \item \textbf{Affinity-based}: Route related tasks to the same agent to exploit cached context.
\end{itemize}

\subsection{Version Management and Compatibility}
\label{version-management-and-compatibility}

Agent Cards include a \texttt{version} field. Orchestrators should specify minimum version requirements and handle graceful degradation when only older versions are available. Semantic versioning~\cite{preston2024semver} (\texttt{MAJOR.MINOR.PATCH}) is recommended: breaking interface changes increment \texttt{MAJOR}, new capabilities increment \texttt{MINOR}.

\begin{warningbox}[Version Skew in Long-Running Systems]
In production multi-agent systems, different agents may be updated at different times, creating version skew. An orchestrator compiled against Agent Card v2.1 may encounter agents still running v1.3. Always implement backward-compatible message handling and test cross-version scenarios explicitly.
\end{warningbox}

\section{Message Formats and Schemas}
\label{sec:a2a:messages}

\subsection{Structured vs.~Unstructured Messages}
\label{structured-vs.-unstructured-messages}

A2A supports a spectrum from fully unstructured (plain text) to fully structured (typed JSON schemas). The right choice depends on the agents involved:

\begin{table}[ht!]
\centering
\caption{Structured vs.~unstructured A2A message trade-offs.}
\begin{tabular}{@{}lp{5cm}p{8cm}@{}}
\toprule
\textbf{Message Type} & \textbf{Advantages} & \textbf{Disadvantages} \\
\midrule
Plain text & Flexible, human-readable, easy to generate & Hard to parse reliably, no schema validation \\
Structured JSON & Machine-parseable, validatable, typed & Requires schema agreement, less flexible \\
Hybrid (text + data) & Human-readable intent + machine-parseable payload & More complex to construct and parse \\
\bottomrule
\end{tabular}
\end{table}

\subsection{Multi-Modal Messages}
\label{multi-modal-messages}

A2A messages are structured as a \textbf{role} (\texttt{user} or \texttt{agent}) plus a list of typed \textbf{parts}:

\begin{table}[ht!]
\centering
\caption{A2A message part types (wire format uses \texttt{"type": "text"|"file"|"data"}).}
\begin{tabular}{@{}lp{5cm}p{8cm}@{}}
\toprule
\textbf{Part Type} & \textbf{Fields} & \textbf{Use Case} \\
\midrule
\texttt{TextPart} & \texttt{text: string} & Natural language instructions, responses \\
\texttt{FilePart} & \texttt{mimeType}, \texttt{uri} or \texttt{bytes} & Documents, images, audio, code files \\
\texttt{DataPart} & \texttt{data: object} & Structured JSON (tool results, schemas) \\
\bottomrule
\end{tabular}
\end{table}

Modern agents increasingly work with non-text modalities. A2A’s \texttt{FilePart} supports any MIME type, enabling rich multi-modal workflows:

\begin{examplebox}[Multi-Modal A2A Message: Data Analysis]
\begin{lstlisting}[style=pythonstyle]
# A message combining text instructions with a data payload and a file
message = {
    "role": "user",
    "parts": [
        {
            "type": "text",
            "text": "Analyze the attached CSV and the schema below. "
                    "Identify anomalies and produce a summary report."
        },
        {
            "type": "file",
            "mimeType": "text/csv",
            "uri": "https://storage.example.com/data/sales_q3.csv"
        },
        {
            "type": "data",
            "data": {
                "schema": {
                    "columns": ["date", "region", "product", "revenue", "units"],
                    "types":   ["date", "string", "string", "float", "int"]
                },
                "expectedRowCount": 15000,
                "anomalyThreshold": 3.0  # z-score threshold
            }
        }
    ]
}
\end{lstlisting}
\end{examplebox}

\begin{examplebox}[Multi-Modal A2A Message: Image Analysis]
\begin{lstlisting}[style=pythonstyle]
# Multi-modal message: text + image + structured data
multimodal_message = {
    "role": "user",
    "parts": [
        {"type": "text",
         "text": "Describe what is wrong with this chart and suggest fixes."},
        {"type": "file",
         "mimeType": "image/png",
         "bytes": base64.b64encode(chart_image_bytes).decode()},
        {"type": "data",
         "data": {
             "chartType": "bar",
             "dataSource": "Q3 Revenue by Region",
             "knownIssues": ["y-axis does not start at zero",
                             "missing error bars"]
         }}
    ]
}
\end{lstlisting}
\end{examplebox}

\subsection{Context Passing: What to Share vs.~What to Keep Private}
\label{context-passing-what-to-share-vs.-what-to-keep-private}

A critical design decision in multi-agent systems is \emph{context scoping}: how much of the conversation history and internal state to pass to a sub-agent.

\begin{keybox}[Context Scoping Principles]
Minimal Context

Pass only what the sub-agent needs to complete its task. Reduces token usage, latency, and the risk of leaking sensitive information.

Summarized Context

Instead of passing raw conversation history, pass a structured summary: goals, constraints, decisions made, and relevant facts.

Private State

Internal reasoning, intermediate drafts, and user PII should generally \emph{not} be forwarded to sub-agents unless explicitly required.

Correlation IDs

Always pass a \texttt{correlationId} so that sub-agent actions can be traced back to the originating workflow in logs and audit trails.
\end{keybox}

\subsection{Conversation Threading and Correlation IDs}
\label{conversation-threading-and-correlation-ids}

In complex workflows, many tasks may be in flight simultaneously. \textbf{Correlation IDs} link related tasks across agents:

\begin{lstlisting}[style=pythonstyle]
import uuid

class WorkflowContext:
    """Carries correlation metadata through a multi-agent workflow."""

    def __init__(self, workflow_id: str | None = None):
        self.workflow_id = workflow_id or str(uuid.uuid4())
        self.span_id = str(uuid.uuid4())
        self.parent_span_id: str | None = None

    def child_context(self) -> "WorkflowContext":
        """Create a child context for a sub-task."""
        child = WorkflowContext(workflow_id=self.workflow_id)
        child.parent_span_id = self.span_id
        return child

    def to_metadata(self) -> dict:
        return {
            "x-workflow-id": self.workflow_id,
            "x-span-id": self.span_id,
            "x-parent-span-id": self.parent_span_id
        }

# Usage: attach to every A2A task submission
ctx = WorkflowContext()
task = await client.send_task(
    message=message,
    metadata=ctx.to_metadata()
)
# Sub-tasks use child contexts for tracing
sub_ctx = ctx.child_context()
\end{lstlisting}

\section{Coordination Protocols}
\label{sec:a2a:coordination}

Beyond point-to-point communication, multi-agent systems benefit from higher-level \textbf{coordination protocols}---structured interaction patterns that enable collective decision-making and problem-solving.

\subsection{Contract Net Protocol}
\label{contract-net-protocol}

The \textbf{Contract Net Protocol (CNP)}~\cite{smith1980contract} is a classic multi-agent coordination mechanism adapted for LLM-based systems:

\begin{enumerate}
  \item \textbf{Announcement}: The manager agent broadcasts a task announcement to all potential contractor agents, including task requirements and evaluation criteria.
  \item \textbf{Bidding}: Contractor agents evaluate the task against their capabilities and submit bids containing estimated completion time, confidence, and resource requirements.
  \item \textbf{Award}: The manager selects the winning bid (or multiple bids for parallel subtasks) and awards the contract.
  \item \textbf{Execution and Reporting}: The contractor executes the task and reports results back to the manager.
\end{enumerate}

\begin{examplebox}[Contract Net Protocol Implementation]
\begin{lstlisting}[style=pythonstyle]
import dataclasses

class ContractNetManager:
    """Implements the Contract Net Protocol for task allocation."""

    async def allocate_task(self, task: Task,
                            candidate_agents: list[AgentCard]) -> AgentCard:
        # Phase 1: Announce task to all candidates
        announcement = {
            "type": "task-announcement",
            "task": dataclasses.asdict(task),
            "deadline": (datetime.now(timezone.utc) + timedelta(seconds=10)).isoformat(),
            "evaluationCriteria": ["confidence", "estimatedTime", "cost"]
        }

        # Phase 2: Collect bids
        bids = await asyncio.gather(*[
            self._request_bid(agent, announcement)
            for agent in candidate_agents
        ], return_exceptions=True)

        valid_bids = [(agent, bid) for agent, bid in zip(candidate_agents, bids)
                      if not isinstance(bid, Exception) and bid is not None]

        if not valid_bids:
            raise RuntimeError(f"No agents bid on task {task.id}")

        # Phase 3: Award to best bidder (highest confidence, lowest time)
        def score_bid(agent_bid):
            _, bid = agent_bid
            return bid["confidence"] - 0.1 * bid["estimatedSeconds"]

        winner_agent, winning_bid = max(valid_bids, key=score_bid)

        # Notify winner and losers
        await self._award_contract(winner_agent, task)
        await asyncio.gather(*[
            self._reject_bid(agent, task.id)
            for agent, _ in valid_bids if agent != winner_agent
        ])

        return winner_agent

    async def _request_bid(self, agent: AgentCard,
                           announcement: dict) -> dict | None:
        """Ask an agent to bid on a task."""
        try:
            result = await self.client.send_task(
                agent_url=agent.url,
                message={"role": "user",
                         "parts": [{"type": "data", "data": announcement}]}
            )
            return result.artifacts[0].parts[0]["data"]
        except Exception:
            return None
\end{lstlisting}
\end{examplebox}

\subsection{Blackboard Systems}
\label{blackboard-systems}

A \textbf{blackboard system}~\cite{hayes1985blackboard} provides a shared workspace (the ``blackboard'') where agents post partial solutions, observations, and hypotheses. Other agents monitor the blackboard and contribute when they can add value---an \emph{opportunistic} problem-solving approach.

Blackboard systems are well-suited to problems where the solution path is not known in advance and different agents may contribute at different stages---such as scientific hypothesis generation, complex debugging, or multi-source intelligence analysis.

\subsection{Consensus Protocols}
\label{consensus-protocols}

When multiple agents must agree on a decision (e.g., which plan to execute, whether a result is correct), \textbf{consensus protocols} provide structured voting mechanisms:

Simple Majority Voting

Each agent votes; the option with $> 50\%$ of votes wins. Fast but vulnerable to correlated errors if agents share the same base model.

Weighted Voting

Votes are weighted by agent confidence or historical accuracy. More robust but requires calibrated confidence estimates.

Quorum-Based

A decision requires agreement from at least $k$ of $n$ agents. Provides fault tolerance: up to $n-k$ agents can fail or disagree without blocking.

Delphi Method

Agents vote, see anonymized results, revise their votes, and repeat until convergence. Reduces anchoring bias and encourages genuine deliberation.

\begin{lstlisting}[style=pythonstyle]
async def quorum_vote(agents: list[AgentCard], question: str,
                      options: list[str], quorum: int) -> str | None:
    """Run a quorum vote across agents. Returns winning option or None."""
    votes = await asyncio.gather(*[
        ask_agent_to_vote(agent, question, options)
        for agent in agents
    ])

    counts: dict[str, int] = {}
    for vote in votes:
        if vote in options:
            counts[vote] = counts.get(vote, 0) + 1

    # Return first option that reaches quorum
    for option, count in sorted(counts.items(), key=lambda x: -x[1]):
        if count >= quorum:
            return option
    return None  # No quorum reached
\end{lstlisting}

\subsection{Leader Election}
\label{leader-election}

In dynamic multi-agent systems, a \textbf{leader} (orchestrator) may need to be elected at runtime---for example, when the original orchestrator fails or when agents self-organize without a pre-assigned coordinator. Classic distributed systems algorithms (Bully, Ring) can be adapted for agent networks, with agents exchanging capability scores or priority tokens to elect the most capable available agent as leader.

\section{A2A vs.~MCP: Complementary Protocols}
\label{sec:a2a:vs_mcp}

A common source of confusion is the relationship between A2A and the \textbf{Model Context Protocol (MCP)}~\cite{anthropic-mcp-2024}. These protocols are \emph{complementary}, not competing:

\begin{keybox}[The Core Distinction]
\begin{itemize}
  \item \textbf{MCP} is the \emph{vertical} protocol: it extends an agent downward into the world of databases, APIs, file systems, and code executors. Only the agent reasons; MCP endpoints are deterministic services.
  \item \textbf{A2A} is the \emph{horizontal} protocol: it links one reasoning agent to another. Both sides are intelligent actors capable of reasoning, planning, and tool use.
\end{itemize}
\end{keybox}

\begin{tabular}{@{}lp{5cm}p{8cm}@{}}
\toprule
\textbf{Dimension} & \textbf{MCP} & \textbf{A2A} \\
\midrule
\textbf{Participants} & Agent $\leftrightarrow$ Tool/Resource & Agent $\leftrightarrow$ Agent \\
\textbf{Intelligence} & One side (agent) is intelligent & Both sides are intelligent \\
\textbf{Statefulness} & Typically stateless tool calls & Stateful tasks with lifecycle \\
\textbf{Streaming} & Limited (tool results) & First-class SSE streaming \\
\textbf{Discovery} & Tool manifests & Agent Cards \\
\textbf{Auth model} & Server-controlled & Mutual, OAuth 2.0 \\
\textbf{Typical latency} & Milliseconds & Seconds to minutes \\
\textbf{Use case} & ``Search the web'', ``Run SQL'' & ``Delegate to specialist'' \\
\bottomrule
\end{tabular}

\subsection{When to Use Which}
\label{when-to-use-which}

\begin{itemize}
  \item Use \textbf{MCP} when the remote endpoint is a deterministic function: a database query, an API call, a code execution sandbox. The agent controls the interaction entirely.
  \item Use \textbf{A2A} when the remote endpoint needs to \emph{reason} about the request: interpret ambiguous instructions, make judgment calls, use its own tools, or engage in multi-turn dialogue.
  \item Use \textbf{both} in the same system: an orchestrator agent uses A2A to delegate to specialist agents, and each specialist agent uses MCP to access its tools.
\end{itemize}

\subsection{Combined Architecture}
\label{combined-architecture}

In production multi-agent systems, A2A and MCP work together at different layers: \textbf{A2A} handles inter-agent delegation and coordination (horizontal communication between peers), while \textbf{MCP} handles each agent’s connection to its tools and data sources (vertical integration with capabilities). This separation of concerns is key to building scalable agentic architectures.

\begin{figure}[ht!]
\centering
\includegraphics[width=0.85\textwidth]{figures/fig_068_combined-a2a-mcp.png}
\caption{Combined A2A + MCP architecture. The orchestrator delegates to specialist agents via A2A; each agent accesses its tools via MCP servers.}
\label{fig:combined-a2a-mcp}
\end{figure}

\begin{tcolorbox}
\begin{itemize}
  \item \textbf{A2A for delegation}: When an agent needs capabilities it doesn’t have, it delegates to another agent via A2A task messages. Each agent is a self-contained service with its own Agent Card.
  \item \textbf{MCP for tool access}: Each agent connects to its tools through MCP servers. This means tools are never exposed directly to other agents --- only through the owning agent’s interface.
  \item \textbf{Separation of trust boundaries}: The orchestrator trusts specialist agents (verified via A2A authentication). Each specialist trusts its own MCP servers (local or authenticated). No transitive tool access.
  \item \textbf{Independent scaling}: Code-heavy workloads can scale CodeAgent instances; data workloads scale DataAgent. The orchestrator remains lightweight.
\end{itemize}
\end{tcolorbox}

\section{Security and Trust in Multi-Agent Systems}
\label{sec:a2a:security}

Multi-agent systems introduce unique security challenges. When Agent A delegates to Agent B, which delegates to Agent C, the chain of trust must be carefully managed.

\subsection{Agent Identity Verification}
\label{agent-identity-verification}

Each agent must have a verifiable identity. Options include:

\begin{itemize}
  \item \textbf{JWT tokens}~\cite{rfc7519} signed by a trusted identity provider, carrying the agent’s ID, issuer, and expiry. Verified by the receiving agent using the provider’s public key.
  \item \textbf{mTLS certificates}~\cite{rfc8705} issued by an internal CA, providing both authentication and transport encryption.
  \item \textbf{Decentralized identifiers (DIDs)}~\cite{w3c-did-2022} for cross-organization scenarios where no single trusted authority exists.
\end{itemize}

\subsection{Message Integrity and Encryption}
\label{message-integrity-and-encryption}

\begin{itemize}
  \item All A2A communication should occur over \textbf{TLS 1.3}~\cite{rfc8446} to prevent eavesdropping and man-in-the-middle attacks.
  \item For sensitive payloads, \textbf{end-to-end encryption} (e.g., JWE) ensures that intermediate infrastructure (load balancers, proxies) cannot read message content.
  \item \textbf{Message signing} (JWS) provides non-repudiation: the receiving agent can prove that a specific message came from a specific sender.
\end{itemize}

\subsection{Authorization Scopes}
\label{authorization-scopes}

Not every agent should be able to ask every other agent to do anything. OAuth~2.0 authorization scopes~\cite{rfc6749} define the boundaries:

\begin{lstlisting}[style=pythonstyle]
# Example OAuth 2.0 scopes for a DataAgent
SCOPES = {
    "data:read":        "Read data from connected databases",
    "data:write":       "Write or modify data in connected databases",
    "data:export":      "Export data to external systems",
    "analysis:run":     "Execute statistical analyses",
    "analysis:schedule":"Schedule recurring analyses",
    "admin:config":     "Modify agent configuration"
}

# A ReportingAgent might hold only: data:read, analysis:run
# An ETL pipeline agent might hold: data:read, data:write, data:export
# Only a human admin holds: admin:config

class A2AServer:
    def verify_authorization(self, token: str, required_scope: str) -> bool:
        """Verify that the calling agent holds the required scope."""
        claims = jwt.decode(token, self.public_key, algorithms=["RS256"])
        granted_scopes = claims.get("scope", "").split()
        if required_scope not in granted_scopes:
            raise PermissionError(
                f"Caller lacks required scope '{required_scope}'. "
                f"Granted: {granted_scopes}"
            )
        return True
\end{lstlisting}

\subsection{Audit Trails and Accountability}
\label{audit-trails-and-accountability}

\begin{warningbox}[The Accountability Gap]
In a chain of agent delegations, it can become unclear \emph{who} is responsible for an action. If Agent A asks Agent B to delete a file, and Agent B does so, who is accountable? Every A2A interaction must be logged with: the calling agent’s identity, the task description, the authorization token used, the timestamp, and the outcome. This audit trail is essential for incident response, compliance, and debugging.
\end{warningbox}

Every A2A server should emit structured audit logs:

\begin{lstlisting}[style=pythonstyle]
@dataclass
class A2AAuditEvent:
    timestamp: str          # ISO 8601
    workflow_id: str        # Correlation ID for the top-level workflow
    span_id: str            # This task's span
    parent_span_id: str     # Calling task's span (for delegation chains)
    caller_agent_id: str    # Verified identity of the calling agent
    callee_agent_id: str    # This agent's identity
    task_id: str
    skill_invoked: str
    authorization_scopes: list[str]
    outcome: str            # "completed" | "failed" | "rejected"
    duration_ms: int
    error_code: str | None
\end{lstlisting}

\section{Implementation Example: Multi-Agent Research Workflow}
\label{sec:a2a:implementation}

The following example demonstrates a complete multi-agent research workflow using A2A: an \texttt{OrchestratorAgent} decomposes a research question, delegates to specialist agents, and synthesizes their results.

\begin{lstlisting}[style=pythonstyle]
"""
Multi-agent research workflow using A2A protocol.
Demonstrates: Agent Cards, A2A client/server, task lifecycle,
multi-turn interaction, and agent handoffs.
"""

import asyncio
import json
import uuid
from collections.abc import AsyncIterator
from datetime import datetime, timedelta, timezone

import httpx
from fastapi import FastAPI, HTTPException, Request
from fastapi.responses import StreamingResponse
from pydantic import BaseModel, Field

# -- Data Models --------------------------------------------------------------

class Part(BaseModel):
    type: str           # "text" | "file" | "data"
    text: str | None = None
    data: dict | None = None
    mimeType: str | None = None
    uri: str | None = None

class Message(BaseModel):
    role: str           # "user" | "agent"
    parts: list[Part]

class TaskStatus(BaseModel):
    state: str          # submitted | working | input-required | completed | failed
    message: str | None = None
    timestamp: str = Field(
        default_factory=lambda: datetime.now(timezone.utc).isoformat()
    )

class Artifact(BaseModel):
    parts: list[Part]
    index: int = 0
    append: bool = False
    lastChunk: bool = True

class Task(BaseModel):
    id: str
    status: TaskStatus
    messages: list[Message] = []
    artifacts: list[Artifact] = []
    metadata: dict = {}

# -- A2A Client (HTTP/REST binding) --------------------------------------------
# Note: A2A v1.0 defines three protocol bindings: JSON-RPC 2.0, gRPC, and
# HTTP+JSON/REST. This example uses the REST binding for readability.

class A2AClient:
    """Client for sending tasks to A2A-compliant agents."""

    def __init__(self, agent_url: str, auth_token: str):
        self.agent_url = agent_url.rstrip("/")
        self.headers = {
            "Authorization": f"Bearer {auth_token}",
            "Content-Type": "application/json"
        }

    async def get_agent_card(self) -> dict:
        """Fetch the agent's capability card."""
        async with httpx.AsyncClient() as client:
            resp = await client.get(
                f"{self.agent_url}/.well-known/agent.json",
                headers=self.headers
            )
            resp.raise_for_status()
            return resp.json()

    async def send_task(self, message: Message,
                        task_id: str | None = None,
                        metadata: dict | None = None) -> Task:
        """Submit a task and return the initial task object."""
        payload = {
            "id": task_id or str(uuid.uuid4()),
            "message": message.model_dump(),
            "metadata": metadata or {}
        }
        async with httpx.AsyncClient() as client:
            resp = await client.post(
                f"{self.agent_url}/tasks/send",
                json=payload,
                headers=self.headers,
                timeout=30.0
            )
            resp.raise_for_status()
            return Task(**resp.json())

    async def stream_task(self, message: Message,
                          metadata: dict | None = None) -> AsyncIterator[dict]:
        """Submit a task and stream SSE events."""
        payload = {
            "id": str(uuid.uuid4()),
            "message": message.model_dump(),
            "metadata": metadata or {}
        }
        async with httpx.AsyncClient() as client:
            async with client.stream(
                "POST",
                f"{self.agent_url}/tasks/sendSubscribe",
                json=payload,
                headers={**self.headers, "Accept": "text/event-stream"},
                timeout=300.0
            ) as response:
                async for line in response.aiter_lines():
                    if line.startswith("data: "):
                        event_data = json.loads(line[6:])
                        yield event_data
                        if event_data.get("final"):
                            break

    async def get_task(self, task_id: str) -> Task:
        """Poll for task status."""
        async with httpx.AsyncClient() as client:
            resp = await client.get(
                f"{self.agent_url}/tasks/{task_id}",
                headers=self.headers
            )
            resp.raise_for_status()
            return Task(**resp.json())

    async def wait_for_completion(self, task: Task,
                                  poll_interval: float = 2.0) -> Task:
        """Poll until task reaches a terminal state."""
        terminal_states = {"completed", "failed", "canceled"}
        while task.status.state not in terminal_states:
            await asyncio.sleep(poll_interval)
            task = await self.get_task(task.id)
        return task

# -- A2A Server (FastAPI) -----------------------------------------------------

class ResearchAgent:
    """
    A specialist research agent that searches literature and
    summarizes findings on a given topic.
    """

    AGENT_CARD = {
        "name": "ResearchAgent",
        "description": "Searches academic literature and synthesizes research findings.",
        "url": "https://research-agent.example.com/a2a",
        "version": "1.0.0",
        "capabilities": {
            "streaming": True,
            "pushNotifications": False,
            "stateTransitionHistory": True
        },
        "authentication": {"schemes": ["Bearer"]},
        "skills": [{
            "id": "literature-search",
            "name": "Literature Search",
            "description": "Search and summarize academic papers on a topic.",
            "tags": ["research", "literature", "academic", "papers"],
            "examples": [
                "Summarize recent papers on transformer attention mechanisms.",
                "What does the literature say about RLHF for code generation?"
            ],
            "inputModes": ["text"],
            "outputModes": ["text", "data"]
        }]
    }

    def __init__(self):
        self.tasks: dict[str, Task] = {}
        self.app = FastAPI(title="ResearchAgent A2A Server")
        self._register_routes()

    def _register_routes(self):
        @self.app.get("/.well-known/agent.json")
        async def agent_card():
            return self.AGENT_CARD

        @self.app.post("/tasks/send")
        async def send_task(request: Request):
            body = await request.json()
            task = await self._create_and_run_task(body)
            return task.model_dump()

        @self.app.post("/tasks/sendSubscribe")
        async def send_subscribe(request: Request):
            body = await request.json()
            return StreamingResponse(
                self._stream_task(body),
                media_type="text/event-stream"
            )

        @self.app.get("/tasks/{task_id}")
        async def get_task(task_id: str):
            if task_id not in self.tasks:
                raise HTTPException(status_code=404, detail="Task not found")
            return self.tasks[task_id].model_dump()

    async def _create_and_run_task(self, body: dict) -> Task:
        task_id = body.get("id", str(uuid.uuid4()))
        message = Message(**body["message"])

        task = Task(
            id=task_id,
            status=TaskStatus(state="submitted"),
            messages=[message],
            metadata=body.get("metadata", {})
        )
        self.tasks[task_id] = task

        # Run asynchronously
        asyncio.create_task(self._execute_task(task_id))
        return task

    async def _execute_task(self, task_id: str):
        task = self.tasks[task_id]
        task.status = TaskStatus(state="working")

        try:
            # Extract the research question from the message
            question = task.messages[0].parts[0].text

            # Simulate literature search (replace with real search tool)
            await asyncio.sleep(1)  # Simulated latency
            findings = await self._search_literature(question)

            # Produce artifact
            task.artifacts = [Artifact(parts=[
                Part(type="text", text=findings["summary"]),
                Part(type="data", data={"papers": findings["papers"],
                                        "query": question})
            ])]
            task.status = TaskStatus(state="completed")

        except Exception as e:
            task.status = TaskStatus(state="failed", message=str(e))

        self.tasks[task_id] = task

    async def _search_literature(self, question: str) -> dict:
        """Placeholder: in production, calls a real search API."""
        return {
            "summary": f"Based on a search of recent literature regarding "
                       f"'{question}', key findings include: ...",
            "papers": [
                {"title": "Attention Is All You Need", "year": 2017,
                 "relevance": 0.95},
                {"title": "RLHF: Training Language Models to Follow Instructions",
                 "year": 2022, "relevance": 0.88}
            ]
        }

    async def _stream_task(self, body: dict) -> AsyncIterator[str]:
        task = await self._create_and_run_task(body)

        # Stream status updates
        yield f"data: {json.dumps({'id': task.id, 'status': {'state': 'submitted'}, 'final': False})}\n\n"
        yield f"data: {json.dumps({'id': task.id, 'status': {'state': 'working'}, 'final': False})}\n\n"

        # Wait for completion
        while task.status.state not in ("completed", "failed", "canceled"):
            await asyncio.sleep(0.5)
            task = self.tasks[task.id]

        # Stream the artifact
        if task.artifacts:
            for part in task.artifacts[0].parts:
                event = {
                    "id": task.id,
                    "artifact": {
                        "parts": [part.model_dump()],
                        "index": 0,
                        "append": False,
                        "lastChunk": True
                    },
                    "final": False
                }
                yield f"data: {json.dumps(event)}\n\n"

        # Final status
        yield f"data: {json.dumps({'id': task.id, 'status': task.status.model_dump(), 'final': True})}\n\n"

# -- Orchestrator: Multi-Agent Workflow ----------------------------------------

class ResearchOrchestrator:
    """
    Orchestrates a multi-agent research workflow:
    1. Decomposes the research question into sub-questions
    2. Dispatches each sub-question to a ResearchAgent
    3. Synthesizes results into a final report
    """

    def __init__(self, research_agent_url: str, auth_token: str):
        self.research_client = A2AClient(research_agent_url, auth_token)
        self.workflow_id = str(uuid.uuid4())

    async def run(self, research_question: str) -> str:
        print(f"[Orchestrator] Starting workflow {self.workflow_id}")
        print(f"[Orchestrator] Question: {research_question}")

        # Step 1: Decompose into sub-questions
        sub_questions = self._decompose(research_question)
        print(f"[Orchestrator] Decomposed into {len(sub_questions)} sub-questions")

        # Step 2: Dispatch sub-questions in parallel
        tasks = await asyncio.gather(*[
            self.research_client.send_task(
                message=Message(role="user", parts=[Part(type="text", text=q)]),
                metadata={"workflowId": self.workflow_id, "subQuestion": i}
            )
            for i, q in enumerate(sub_questions)
        ])

        # Step 3: Wait for all tasks to complete
        completed_tasks = await asyncio.gather(*[
            self.research_client.wait_for_completion(task)
            for task in tasks
        ])

        # Step 4: Check for failures
        failed = [t for t in completed_tasks if t.status.state == "failed"]
        if failed:
            print(f"[Orchestrator] Warning: {len(failed)} sub-tasks failed")

        # Step 5: Synthesize results
        findings = []
        for task, question in zip(completed_tasks, sub_questions):
            if task.status.state == "completed" and task.artifacts:
                summary = task.artifacts[0].parts[0].text
                findings.append(f"### {question}\n{summary}")

        report = self._synthesize(research_question, findings)
        print(f"[Orchestrator] Workflow complete. Report: {len(report)} chars")
        return report

    def _decompose(self, question: str) -> list[str]:
        """Decompose a complex question into focused sub-questions."""
        # In production: use an LLM to decompose
        return [
            f"What are the foundational methods for: {question}?",
            f"What are the most recent advances in: {question}?",
            f"What are the open challenges and limitations in: {question}?"
        ]

    def _synthesize(self, question: str, findings: list[str]) -> str:
        """Synthesize sub-findings into a coherent report."""
        # In production: use an LLM to synthesize
        sections = "\n\n".join(findings)
        return f"# Research Report: {question}\n\n{sections}"

# -- Entry Point ---------------------------------------------------------------

async def main():
    orchestrator = ResearchOrchestrator(
        research_agent_url="https://research-agent.example.com/a2a",
        auth_token="eyJhbGciOiJSUzI1NiJ9..."
    )
    report = await orchestrator.run(
        "Reinforcement learning from human feedback for large language models"
    )
    print(report)

if __name__ == "__main__":
    asyncio.run(main())
\end{lstlisting}

\section{Summary}
\label{sec:a2a:summary}

\begin{keybox}[Key Takeaways: Agent-to-Agent Communication]
\begin{enumerate}
  \item \textbf{A2A enables specialization at scale}: By routing tasks to specialist agents, multi-agent systems achieve depth and breadth simultaneously. (Chapter~\ref{sec:multi-agent-systems} covers multi-agent architectures in depth.)
  \item \textbf{Google’s A2A Protocol} provides a production-ready, open standard for interoperable agent communication, with Agent Cards, task lifecycle management, SSE streaming, and enterprise authentication.
  \item \textbf{Communication patterns} range from simple request-response to complex negotiation and auction-based allocation---choose based on task complexity and latency needs.
  \item \textbf{A2A and MCP are complementary}: A2A connects agents to agents; MCP connects agents to tools. Most production systems use both.
  \item \textbf{Security is non-negotiable}: Agent identity verification, authorization scopes, and audit trails are essential in any multi-agent deployment.
  \item \textbf{Coordination protocols} (Contract Net, Blackboard, Consensus) provide structured mechanisms for collective decision-making beyond simple delegation.
  \item \textbf{Observability through correlation IDs} is critical for debugging and auditing complex multi-agent workflows spanning many agents and tools.
\end{enumerate}
\end{keybox}

\begin{questionbox}[Open Research Questions in A2A]
\begin{itemize}
  \item How should agents handle \emph{conflicting instructions} from multiple orchestrators in a hierarchy? What conflict resolution mechanisms are most effective?
  \item Can agents \emph{learn} better routing and delegation strategies through experience, rather than relying on static capability declarations?
  \item How do we prevent \emph{prompt injection attacks} where a malicious agent manipulates a downstream agent by embedding adversarial instructions in its messages?
  \item What are the right \emph{privacy boundaries} for context passing---how much conversation history should a sub-agent see, and how do we enforce these boundaries technically?
  \item As agent networks grow to hundreds or thousands of agents, how do we maintain \emph{coherent global state} without creating bottlenecks or consistency violations?
\end{itemize}
\end{questionbox}

